%% =====================================================================
%%  T-16-05  The Summary That Gets "That's Right"
%%  -------------------------------------------------------------------
%%  One idea per file. Core is mandatory. Slots in canonical order.
%%  Max 700 words. No layout commands. No filler. New source material
%%  for this entry goes in sources/<BOOK>/T-16-05.tex -- never by
%%  rewriting this file (docs/RULES.md R0.1).
%% =====================================================================
%% @kind: topic
%% @id: T-16-05
%% @title: The Summary That Gets "That's Right"
%% @subtitle: Felt agreement without conceding the argument
%% @chapter: c16
%% @order: 50
%% @status: stable
%% @sources: B5
%% @updated: 2026-09-19

\topic{T-16-05}{The Summary That Gets ``That's Right''}{Felt agreement without conceding the argument}

\begin{core}
Restate the counterpart's position so completely and accurately --- their words, their reasons, their feelings --- that they answer \emph{that's right}. In that moment they have agreed with you about their own situation, which is a real act of commitment bought without a single concession, and it opens movement that argument cannot.
\end{core}
\begin{mechanism}
A summary combines paraphrase with a closing label. Its power is recognition: most people have never heard their own position stated back to them correctly, and the experience of being accurately heard produces the one answer the source treats as a milestone --- \emph{that's right} --- which is felt as internal agreement, not extracted compliance. The distinction does the diagnostic work: \emph{you're right} is the counterfeit, the phrase people use to end a conversation they have stopped engaging with. Commitment attaches to conclusions the person reaches; a summary makes your framing the occasion for their conclusion.
\end{mechanism}
\begin{conditions}
Strongest after a long or tangled statement of position, with emotional or stubborn counterparts, and any time agreement exists only as counterfeit yeses. Weak against a counterpart whose position is a rote script (nothing true left to reflect), in genuinely simple transactions where the ceremony wastes time, and against a counterpart who recognises therapy-style handling and resents it.
\end{conditions}
\begin{application}
  \item Collect the parts as they talk: their stated position, their reasons, and the emotion under it.
  \item Deliver the whole picture back in one calm passage --- their case, better organised than they made it, ending on their stake.
  \item Close with a label of the feeling underneath it, then stop.
  \item Wait for \emph{that's right}. Do not move before it arrives.
  \item Only then introduce your ask --- as the natural next step inside a world they have just agreed is accurately described.
\end{application}
\begin{feedback}
  \item Landing: \emph{that's right} --- sometimes with visible relief; posture opens and the counterpart starts problem-solving aloud.
  \item Landing: they add the missing piece unprompted, correcting the summary toward their real constraint.
  \item Failing: \emph{you're right} --- the conversation is over in their head; reopen with a mirror (T-07-04) or a label, not with your ask.
  \item Failing: corrections delivered flatly with no engagement. The summary missed the stake; re-run it on what they corrected.
\end{feedback}
\begin{failure}
A summary that is technically complete but emotionally toneless buys nothing --- accuracy without recognition reads as transcription. It can also hand over structure: a precise reflection of their position tells a sharp counterpart you have organised their case better than they have, and some will go quiet rather than be handled. And the milestone can be faked back --- \emph{that's right} can be flattery from a counterpart running the same play.
\end{failure}
\begin{countermeasures}
  \item Distinguish the two phrases in your own mouth. If you are saying \emph{you're right} to end a conversation, notice what you just declined to engage with.
  \item When someone summarises you accurately, you are half-agreed. Ask what they want before accepting the frame their summary implies.
  \item Treat a too-perfect reflection as technique: answer its accuracy, then state your price for continuing --- \emph{that's right, and here is what I need}.
  \item Beware your own relief at being heard; that relief is the moment influence is being bought.
\end{countermeasures}

\sources{B5}
\seesources
\seealso{T-16-02, T-07-04, T-22-02}
